% Section 7.1, from lectures/L07/appendix/a_eeprom.md.

\section{EEPROM}
\label{c7:sec:eeprom}\label{c7:app:a}

\subsection{EEPROM: what it is and why}
\label{c7:a:1}
\begin{itemize}
  \item RAM forgets everything the instant power is removed. EEPROM (Electrically Erasable
    Programmable ROM) doesn't: it's a separate, much smaller memory that keeps its contents across
    a reset or power loss, useful for settings, calibration values, or a counter that should survive
    being unplugged.
  \item The ATmega328P has 1024 bytes of EEPROM, addressed \code{0}--\code{1023}, entirely separate
    from flash (where your program lives) and RAM (where your variables live).
  \item Two costs come with that durability: a single byte write takes a few milliseconds
    (eternity, compared to a RAM write), and each EEPROM cell only tolerates roughly 100,000 write
    cycles before it wears out. Don't write to EEPROM in a tight loop or once per iteration of
    \code{main()}; write only when a value actually needs to persist.
\end{itemize}

\subsection{EEPROM registers}
\label{c7:a:2}
Three registers, all specific to the EEPROM peripheral:
\begin{enumerate}
  \item \textbf{\code{EEAR}} (EEPROM Address Register) holds the address, \code{0}--\code{1023}, to
    read from or write to.
  \item \textbf{\code{EEDR}} (EEPROM Data Register) holds the byte being written, or receives the
    byte just read.
  \item \textbf{\code{EECR}} (EEPROM Control Register) controls and reports on the read/write
    process:
  \begin{itemize}
    \item \code{EERE} (EEPROM Read Enable) starts a read.
    \item \code{EEMPE} (EEPROM Master Write Enable) must be set immediately before \code{EEPE},
      it's a safety interlock, not a data bit.
    \item \code{EEPE} (EEPROM Write Enable) starts a write, and stays set until the write
      completes; a new write must wait for it to clear first.
  \end{itemize}
\end{enumerate}
The datasheet requires \code{EEPE} to be set within four CPU cycles of \code{EEMPE} being set, or
the write silently doesn't happen at all. An interrupt firing in that window is a real risk (a few
cycles is nothing), so a write disables interrupts for the handful of instructions between the
two:

\begin{ccode}
while (EECR & (1U << EEPE)) {} // Wait for any earlier write to finish.

EEAR = addr; // Set destination address (0-1023).
EEDR = byte; // Set data to write (one byte).

cli();                 // Disable interrupts during the write sequence.
EECR |= (1U << EEMPE); // EEMPE must be followed by EEPE within 4 cycles.
EECR |= (1U << EEPE);
sei();                 // Re-enable interrupts once the write sequence is complete.
\end{ccode}

\noindent Reading doesn't have this timing constraint, no write is happening, so there's no
equivalent interlock, just the same ``wait for the peripheral to be free'' pattern already familiar
from \code{ADIF} (Chapter~\ref{c4:ch}) and \code{UDRE0} (Chapter~\ref{c6:ch}):

\begin{ccode}
while (EECR & (1U << EEPE)) {} // Wait for any earlier write to finish.
EEAR = addr;                   // Set the address to read from.
EECR |= (1U << EERE);          // Start read operation.
return EEDR;                   // Return the byte at the specified address.
\end{ccode}

\subsection{\code{eeprom_write_byte()} / \code{eeprom_read_byte()}}
\label{c7:a:3}

\begin{ccode}
#include <stdint.h>

#include <avr/interrupt.h>
#include <avr/io.h>

/**
 * @brief Write a byte to EEPROM at the given address.
 *
 * @param[in] addr EEPROM address to write to (0-1023).
 * @param[in] byte Byte to write.
 */
static void eeprom_write_byte(const uint16_t addr, const uint8_t byte)
{
    // Wait for any earlier write to finish.
    while (EECR & (1U << EEPE)) {}

    // Set destination address and data.
    EEAR = addr;
    EEDR = byte;

    // Perform write operation.
    cli();
    EECR |= (1U << EEMPE);
    EECR |= (1U << EEPE);
    sei();
}

/**
 * @brief Read a byte from EEPROM at the given address.
 *
 * @param[in] addr EEPROM address to read from (0-1023).
 *
 * @return Byte stored at that address.
 */
static uint8_t eeprom_read_byte(const uint16_t addr)
{
    // Wait for any earlier write to finish.
    while (EECR & (1U << EEPE)) {}

    // Set address to read from.
    EEAR = addr;

    // Start read operation.
    EECR |= (1U << EERE);

    // Return the byte stored at the given address.
    return EEDR;
}
\end{ccode}

\subsection{Worked example: persisting LED1's state across resets}
\label{c7:a:4}
Extends \secref{c6:a:6}'s button-toggle-report demo: LED1's state now survives a reset or power
cycle, read back from EEPROM in \code{setup()} and written back out every time it changes. The
circuit is the one of Figure~\ref{c6:fig:serial}, unchanged.

\begin{console}
eeprom_driver_demo/
├── Makefile
├── main.c
├── include/
│   └── driver/
│       ├── eeprom.h
│       └── serial.h
└── source/
    └── driver/
        ├── eeprom.c
        └── serial.c
\end{console}

\noindent The driver files (\file{driver/eeprom.h} and \file{driver/eeprom.c}, both printed in full
below) wrap \secref{c7:a:3}'s \code{eeprom_write_byte()}/\code{eeprom_read_byte()} behind a
buffer-and-length interface, \code{eeprom_write()}/\code{eeprom_read()}, the same shape as
Chapter~\ref{c6:ch}'s \code{serial_write()}. \code{main.c} in full follows them.

\subsubsection{The driver in full}
\cfile{lectures/L07/eeprom_driver_demo/include/driver/eeprom.h}
\needspace{8\baselineskip}
\cfile{lectures/L07/eeprom_driver_demo/source/driver/eeprom.c}

\subsubsection{The program}
\cfile{lectures/L07/eeprom_driver_demo/main.c}

\noindent\code{ISR(PCINT0_vect)} only toggles LED1 and sets \code{led1_event}; the EEPROM write and
the UART report both happen in \code{main()}'s loop instead, once per event, never inside the ISR
itself. That's deliberate, not incidental, and it's good practice generally: both
\code{eeprom_write()} (which can block for several milliseconds waiting on \code{EECR}'s
\code{EEPE} bit, \secref{c7:a:2}) and \code{serial_print()} (which blocks byte-by-byte on
\code{UDRE0}) are slow, blocking calls, and an ISR that spends milliseconds inside either one holds
up every other interrupt for that whole time, including the debounce timer this very feature
depends on. Keep ISRs to the minimum work needed (flip a pin, set a flag) and let \code{main()}'s
loop handle anything that can block, the same reasoning \secref{c6:a:6} gave for never calling
\code{serial_print()} from an ISR.

\subsubsection{Building and running}
Build and flash the same way as \secref{c6:a:6}: on Linux, run \code{make} in the demo's
directory, \file{lectures/L07/eeprom_driver_demo} (\code{make flash} to program the board); on
Windows via Microchip Studio, recreate this folder structure inside the project and add
\file{include} as an include directory. The demo's \file{Makefile} is the one printed in
\secref{c6:sec:build}, with the EEPROM driver added to its source files:

\begin{makecode}
SRC_FILES := main.c \
             source/driver/eeprom.c \
             source/driver/serial.c \
\end{makecode}
